\documentclass[11pt,a4paper]{article}

\usepackage{../shared/cathedral-preamble}
\usepackage{setspace}
\onehalfspacing

\title{%
  167 Years of Attempts:\\
  The Cathedral in the History of the Riemann Hypothesis%
}
\author{
  Jason Robert Gochanour
}
\date{April 20, 2026}

\begin{document}

\maketitle

\begin{abstract}
We place the Cathedral project in historical context by tracing
the 167-year arc of approaches to the Riemann Hypothesis, from
Riemann's 1859 memoir through the Hilbert--P\'olya conjecture,
the Montgomery--Odlyzko law, the Nyman--Beurling criterion, and
the B\'aez-Duarte strengthening, to the present machine-verified
reduction. We identify three historical turning points where the
approach changed paradigm: (1)~the shift from direct zero-finding
to spectral methods (1914--1973), (2)~the shift from operator
theory to $L^2$ approximation (1950--2003), and (3)~the shift
from informal proof to machine verification (2017--2026). The
Cathedral sits at the confluence of the second and third shifts,
combining B\'aez-Duarte's $L^2$ criterion with Lean~4
verification. We argue that the project's significance lies not
in resolving the conjecture but in \emph{formalizing the map}: for
the first time, the complete logical structure of an RH equivalence
has been machine-checked, with every assumption named, every
deduction verified, and every failed approach preserved.
\end{abstract}

\tableofcontents

% ================================================================
\section{Act I: The Conjecture (1859--1914)}
% ================================================================

\subsection{Riemann's Eight Pages}

In November 1859, Bernhard Riemann presented an eight-page paper
to the Berlin Academy: ``\"Ueber die Anzahl der Primzahlen unter
einer gegebenen Gr\"osse'' (On the number of primes less than a
given magnitude). In it, he introduced the analytic continuation
of the zeta function to the entire complex plane, proved the
functional equation $\xi(s) = \xi(1-s)$, and made a remark that
has haunted mathematics ever since:

\begin{quote}
\emph{``One would of course like to have a rigorous proof of this,
but I have put aside the search for such a proof after some
fleeting vain attempts, because it is not necessary for the
immediate objective of my investigation.''}
\end{quote}

The ``this'' was the observation that all non-trivial zeros of
$\zeta(s)$ have real part $\frac{1}{2}$.

Riemann published one paper on number theory in his life. He died
in 1866, at age 39. The conjecture survived him by 167 years
and counting.

\subsection{Early Attempts}

The first systematic studies of zeta zeros were computational:

\begin{itemize}
  \item \textbf{Gram (1903)}: Computed the first 15 zeros on the
    critical line. Introduced the Gram points $g_n$ where
    $\vartheta(g_n) = n\pi$.
  \item \textbf{Backlund (1914)}: Verified RH for the first 200
    zeros.
  \item \textbf{Hardy (1914)}: Proved that infinitely many zeros
    lie on the critical line, without proving that \emph{all} do.
\end{itemize}

These early efforts established a pattern that would persist for
a century: computational verification of RH for increasing numbers
of zeros, accompanied by partial theoretical progress (``many zeros
are on the line'') but no complete proof.

% ================================================================
\section{Act II: The Spectral Turn (1914--1987)}
% ================================================================

\subsection{Hilbert and P\'olya}

Around 1914, David Hilbert and George P\'olya independently
suggested that the zeta zeros might be eigenvalues of a
self-adjoint operator $T$ on some Hilbert space: $T\psi_n =
\rho_n \psi_n$ where $\rho_n = \frac{1}{2} + i\gamma_n$. If $T$
is self-adjoint, its eigenvalues are real, so $\gamma_n \in \mathbb{R}$
and RH follows.

This was a paradigm shift: instead of studying the zeta function
directly, study the \emph{operator} whose spectrum it encodes.
The search for the ``Hilbert--P\'olya operator'' became a
parallel track to direct analytical approaches.

\subsection{Selberg's Contribution}

In 1942, Atle Selberg proved that a positive proportion of zeros
lie on the critical line. His methods---the Selberg sieve and the
Selberg trace formula---introduced the tools that would eventually
connect number theory to spectral theory and random matrix theory.

The Cathedral's log-cutoff witness $v_k = -\mu(k)(1 - \ln k/\ln N)$
is a Selberg-type sieve weight. The $L^2$ variational principle
independently rediscovers these weights as optimal---a connection
that Selberg himself might have appreciated.

\subsection{Montgomery--Odlyzko: The GUE Connection}

In 1973, Hugh Montgomery discovered that the pair correlation of
zeta zeros matches the Gaussian Unitary Ensemble (GUE) distribution
from random matrix theory. In 1987, Andrew Odlyzko computed
$10^{20}$ zeros and confirmed Montgomery's prediction to remarkable
precision.

This discovery connected the Riemann Hypothesis to quantum physics:
the same statistical patterns appear in the energy levels of heavy
nuclei and the eigenvalues of random matrices. The Hilbert--P\'olya
conjecture went from speculation to a concrete statistical prediction.

\textbf{Connection to the Cathedral}: The Gram matrix $G_N$ in the
Cathedral is a deterministic number-theoretic matrix whose eigenvalue
distribution can be compared with GUE predictions. The numerical
experiments show $\lambda_{\min}(G_N) \sim C/\ln N$, a decay rate
consistent with the soft edge of the Tracy--Widom distribution.

% ================================================================
\section{Act III: The $L^2$ Revolution (1950--2003)}
% ================================================================

\subsection{Nyman (1950)}

Bertil Nyman's 1950 Uppsala thesis established the first function-
space criterion for RH: the Riemann Hypothesis is true if and only
if the constant function $1$ lies in the $L^2(0,1)$ closure of the
span of dilated fractional parts $\{\theta/x\}$, $0 < \theta \leq 1$.

This was the second paradigm shift: instead of studying the zeta
function or its operator, study an \emph{approximation problem}
in a Hilbert space. RH becomes a question about whether certain
functions can approximate a target.

\subsection{Beurling (1955)}

Arne Beurling simplified Nyman's criterion and gave a cleaner
formulation. The Nyman--Beurling theorem states:

\[
  \RH \;\iff\;
  \overline{\operatorname{span}}\{x \mapsto \{\theta/x\} :
  0 < \theta \leq 1\} = L^2(0,1).
\]

\subsection{B\'aez-Duarte (2003)}

Luis B\'aez-Duarte made the crucial strengthening: one can restrict
the dilations to integer reciprocals $\theta = 1/k$, $k = 2, 3,
\ldots$ This restriction is not trivial---it converts the
continuous parameter $\theta$ into a discrete arithmetic parameter,
connecting the approximation problem to the multiplicative structure
of the integers.

The B\'aez-Duarte distance $d_N^2 = \inf_v \int_0^1 |1 - \sum v_k
\{1/(kx)\}|^2\,dx$ is the object the Cathedral studies. The entire
project reduces to understanding this single quantity.

\subsection{Vasyunin and Balazard--Saias}

Vasyunin (1996) showed that the Gram matrix entries $G(j,k) =
\int_0^1 \{1/(jx)\}\{1/(kx)\}\,dx$ have a closed-form expression
involving logarithms, GCDs, and cotangent sums. Balazard and Saias
(2000) established the asymptotic rate $d_N^2 \sim c/\ln N$.

The Cathedral formalizes both of these results, though the crown
theorem ultimately bypasses Vasyunin's cotangent formula via the
Parseval Bridge.

% ================================================================
\section{Act IV: The Verification Turn (2017--2026)}
% ================================================================

\subsection{The Rise of Formal Verification}

The third paradigm shift is methodological rather than mathematical:
the application of machine verification to mathematical proof.

Key milestones:
\begin{itemize}
  \item \textbf{Hales (2017)}: Formal verification of the Kepler
    Conjecture (Flyspeck project), 6 years of effort, HOL Light.
  \item \textbf{Scholze--Buzzard (2021)}: Formalization of
    liquid tensor experiments in Lean, establishing Lean as viable
    for frontier mathematics.
  \item \textbf{Tao et al. (2023)}: Formalization of the
    Polynomial Freiman--Ruzsa conjecture in Lean~4/Blueprint.
  \item \textbf{AlphaProof (2024)}: Google DeepMind's AI achieves
    IMO silver medal level using Lean~4.
\end{itemize}

\subsection{The Cathedral (2026)}

The Cathedral is, to our knowledge, the first machine-verified
formalization of a complete RH equivalence. It differs from prior
work in several respects:

\begin{center}
\small
\begin{tabular}{@{}llccc@{}}
\toprule
\textbf{Project} & \textbf{Result} & \textbf{Prover} &
\textbf{Team} & \textbf{Duration} \\
\midrule
Flyspeck & Kepler conjecture & HOL Light & $\sim$20 & 6 years \\
LTE & Liquid tensors & Lean~3 & $\sim$10 & 6 months \\
PFR & Polynomial FR & Lean~4 & $\sim$15 & 3 months \\
Cathedral & RH equivalence & Lean~4 & 1 + AI & 23 days \\
\bottomrule
\end{tabular}
\end{center}

The comparison is not entirely fair---the Cathedral does not prove
RH, while Flyspeck proves Kepler, and the mathematical content
differs in depth. But the timeline is striking: one person with AI
assistance in 23 days, producing 84 files and 641 theorems.

% ================================================================
\section{What the Cathedral Does Differently}
% ================================================================

\subsection{The Parseval Bridge}

Every previous approach to the forward direction
($\RH \Rightarrow d_N^2 \to 0$) goes through the Vasyunin
cotangent formula, requiring a closed-form computation of the
Gram matrix entries. The Cathedral's Parseval Bridge bypasses this
entirely, working in the frequency domain via Plancherel's theorem.

This is historically significant because the Vasyunin formula
involves Dedekind-type sums whose reciprocity properties are
notoriously difficult to formalize. The Cathedral discovered
(empirically) that a candidate reciprocity law was \emph{false}
at $(a,b) = (3,2)$, motivating the switch to the Parseval Bridge.

\subsection{The Rank-1 Mellin Miracle}

The converse direction ($d_N^2 \to 0 \Rightarrow \RH$) was
traditionally proved via Hahn--Banach separation. The Cathedral
replaces this with a direct computation: at a zeta zero $\rho$,
the Mellin transform of the basis functions $\{1/(kx)\}$
factorizes as $\mathcal{M}[h_k](\rho) = 1/(k(\rho-1))$. The
$k$-dependence and $\rho$-dependence separate, making the
obstruction to approximation explicit and elementary.

This is arguably a new proof technique, not just a formalization
of an existing one.

\subsection{The Archive as Artifact}

Previous formalization projects discard failed attempts. The
Cathedral preserves all 96 archived files as a historical record
of the formalization process. This is unprecedented in formal
verification and serves both educational and methodological
purposes.

% ================================================================
\section{The Arc of History}
% ================================================================

\begin{center}
\small
\begin{tabular}{@{}rlp{7cm}@{}}
\toprule
\textbf{Year} & \textbf{Mathematician} & \textbf{Contribution} \\
\midrule
1859 & Riemann & The conjecture \\
1896 & Hadamard, de la Vall\'ee-Poussin & Prime Number Theorem
  (zeta non-vanishing on $\sigma = 1$) \\
1903 & Gram & First zeros computed \\
1914 & Hardy & Infinitely many zeros on the line \\
$\sim$1914 & Hilbert, P\'olya & Spectral conjecture \\
1942 & Selberg & Positive proportion on the line \\
1950 & Nyman & $L^2$ criterion \\
1955 & Beurling & Simplified criterion \\
1973 & Montgomery & GUE pair correlation \\
1987 & Odlyzko & Numerical confirmation of GUE \\
1996 & Vasyunin & Cotangent sum formula for Gram matrix \\
2003 & B\'aez-Duarte & Integer-restricted basis \\
2004 & Gourdon & $10^{13}$ zeros verified \\
2020 & Platt--Trudgian & $3 \times 10^{12}$ zeros verified
  rigorously \\
2026 & \textbf{Cathedral} & \textbf{Machine-verified RH equivalence} \\
\bottomrule
\end{tabular}
\end{center}

The Cathedral does not end this arc. It adds a new kind of entry:
not a partial result about zeros, not a new operator candidate,
but a \emph{certified map} of the logical territory. For the first
time, the complete biconditional $\RH \iff d_N^2 \to 0$ has been
verified by machine, with every assumption exposed.

% ================================================================
\section{What Has Not Been Done}
% ================================================================

An honest historical account must note what the Cathedral does
\emph{not} accomplish:

\begin{enumerate}
  \item It does not prove RH. No approach has.
  \item It does not find a new zero. Computational verification
    of zeros (Gourdon, Platt--Trudgian) is a separate endeavor.
  \item It does not identify the Hilbert--P\'olya operator. The
    spectral interpretation remains conjectural.
  \item It does not improve the zero-free region. The best known
    zero-free region (Korobov--Vinogradov type) is not addressed.
  \item It does not resolve any Millennium Prize Problem. The
    Clay Mathematics Institute requires a proof of RH itself,
    not an equivalence.
\end{enumerate}

What it does is draw the most precise map ever made of what a
proof would require: seven axioms, named and classified, each
encoding a well-understood piece of classical analysis.

% ================================================================
\section{Conclusion: The Long View}
% ================================================================

In 1900, David Hilbert placed the Riemann Hypothesis eighth on
his famous list of 23 problems for the twentieth century. It
was not solved in the twentieth century. At the turn of the
twenty-first, the Clay Mathematics Institute placed it among the
seven Millennium Prize Problems. It has not been solved in the
first quarter of the twenty-first century either.

The Cathedral is not the solution. But it may be part of the
prelude. Every great proof is preceded by a period of
clarification---a period where the community comes to understand
exactly what needs to be proved and exactly what tools are
available. Riemann's 1859 paper was such a clarification for
analytic number theory. The Nyman--Beurling--B\'aez-Duarte
criterion was such a clarification for the $L^2$ approach. The
Cathedral is such a clarification for the verification approach:
here is the map, here are the assumptions, here is what remains.

The question Riemann set aside ``after some fleeting vain attempts''
in 1859 is still open. But the shape of the question has never
been clearer.

\begin{quote}
\emph{167 years of attempts.\\
40 mathematicians' work, assembled and verified.\\
Seven axioms remain.\\
The next chapter belongs to whoever proves them.}
\end{quote}

\begin{thebibliography}{99}

\bibitem{riemann1859}
B.~Riemann, \emph{\"Ueber die Anzahl der Primzahlen unter einer
gegebenen Gr\"osse}, Monatsberichte der Berliner Akademie, 1859.

\bibitem{hardy1914}
G.H.~Hardy, \emph{Sur les z\'eros de la fonction $\zeta(s)$ de
Riemann}, C. R. Acad. Sci. Paris \textbf{158} (1914), 1012--1014.

\bibitem{selberg1942}
A.~Selberg, \emph{On the zeros of Riemann's zeta-function},
Skr. Norske Vid. Akad. Oslo I \textbf{10} (1942), 1--59.

\bibitem{montgomery1973}
H.L.~Montgomery, \emph{The pair correlation of zeros of the zeta
function}, Proc. Symp. Pure Math. \textbf{24} (1973), 181--193.

\bibitem{odlyzko1987}
A.M.~Odlyzko, \emph{On the distribution of spacings between zeros
of the zeta function}, Math. Comp. \textbf{48} (1987), 273--308.

\bibitem{nyman1950}
B.~Nyman, \emph{On the One-Dimensional Translation Group},
Ph.D. Thesis, Uppsala, 1950.

\bibitem{baezduarte2003}
L.~B\'aez-Duarte, \emph{A strengthening of the Nyman--Beurling
criterion}, Atti Accad. Naz. Lincei \textbf{14} (2003), 5--11.

\bibitem{hales2017}
T.~Hales et al., \emph{A Formal Proof of the Kepler Conjecture},
Forum of Mathematics, Pi \textbf{5} (2017), e2.

\end{thebibliography}

\end{document}
